\documentclass{article}
\usepackage{lmodern}
\usepackage{amsmath}
\usepackage{listings}
\usepackage{fullpage}
\usepackage{parskip}
\usepackage{xcolor}
\lstset{
	basicstyle=\ttfamily,
	columns=fullflexible,
	frame=single,
	breaklines=true,
	postbreak=\mbox{\textcolor{red}{$\hookrightarrow$}\space},
}
\begin{document}
\title{Experiment `p\_4\_2\_c\_cumulative\_totals\_array' Results}
\author{\today}
\date{}
\maketitle
\textbf{Experiment outcome: }SUCCESS
\\\textbf{Bad responses: }0
\\\textbf{Responses containing}~\texttt{assume}~\textbf{: }0
\\\textbf{Responses containing}~\texttt{expect}~\textbf{: }0
\\\textbf{Resolution attempts: }1
\\\textbf{Hard fails (resolution): }0
\\\textbf{Soft fails (resolution): }0
\\\textbf{Verification attempts: }1
\section*{Problem Specification}
\textbf{Problem name: }p\_4\_2\_c\_cumulative\_totals\_array
\\\textbf{Natural language statement: }Write a method that takes a sequence of integers and returns cumulative totals. For example, if the input is [1, 7, 2, 9], the method should return [1, 8, 10, 19].
\\\textbf{Method signature: }p\_4\_2\_c\_cumulative\_totals\_array(inputs: array<int>) returns (totals: array<int>)
\subsection*{Ensures}
\begin{itemize}
\item \texttt{totals.Length == inputs.Length}
\item \texttt{forall i :: 0 <= i < inputs.Length ==> totals[i] == (if i == 0 then inputs[0] else totals[i - 1] + inputs[i])}
\end{itemize}
\subsection*{Requires}
\begin{itemize}
\item \texttt{0 < inputs.Length}
\end{itemize}
\clearpage
\section*{GenAI interactions}
Below you will find all interactions between the `user' (program) and the `assistant' (OpenAI).
\subsection*{Program $\rightarrow$ GenAI}
\begin{lstlisting}
You are given the following task to perform in Dafny:

Write a method that takes a sequence of integers and returns cumulative totals. For example, if the input is [1, 7, 2, 9], the method should return [1, 8, 10, 19].

The signature should be:

p_4_2_c_cumulative_totals_array(inputs: array<int>) returns (totals: array<int>)

The method should respect the following contract:

requires 0 < inputs.Length, ensures totals.Length == inputs.Length, ensures forall i :: 0 <= i < inputs.Length ==> totals[i] == (if i == 0 then inputs[0] else totals[i - 1] + inputs[i])

Produce and show only the Dafny body of this method, including the curly braces that surround it. Do not show the signature nor contract. You must not use 'assume' nor 'expect' anywhere in your code.
\end{lstlisting}
\subsection*{GenAI $\rightarrow$ Program}
\textbf{System fingerprint: }null
\\\textbf{ID: }chatcmpl-Dme8v8f95SahOYIM4e8KSAmwXgf0Q
\\\textbf{Model: }gpt-5.5-2026-04-23
\\\textbf{Created at: }1780486425
\\\textbf{Finish reason: }stop
\\\textbf{Usage: }428 tokens in, and 660 tokens out
\begin{lstlisting}
{
  totals := new int[inputs.Length];
  totals[0] := inputs[0];

  var i := 1;
  while i < inputs.Length
    invariant 1 <= i <= inputs.Length
    invariant totals.Length == inputs.Length
    invariant totals[0] == inputs[0]
    invariant forall j :: 1 <= j < i ==> totals[j] == totals[j - 1] + inputs[j]
  {
    totals[i] := totals[i - 1] + inputs[i];
    i := i + 1;
  }
}
\end{lstlisting}
\section*{Final Program}
\begin{lstlisting}
method p_4_2_c_cumulative_totals_array(inputs: array<int>) returns (totals: array<int>)
	requires 0 < inputs.Length
	ensures totals.Length == inputs.Length
	ensures forall i :: 0 <= i < inputs.Length ==> totals[i] == (if i == 0 then inputs[0] else totals[i - 1] + inputs[i])
{
  totals := new int[inputs.Length];
  totals[0] := inputs[0];

  var i := 1;
  while i < inputs.Length
    invariant 1 <= i <= inputs.Length
    invariant totals.Length == inputs.Length
    invariant totals[0] == inputs[0]
    invariant forall j :: 1 <= j < i ==> totals[j] == totals[j - 1] + inputs[j]
  {
    totals[i] := totals[i - 1] + inputs[i];
    i := i + 1;
  }
}
\end{lstlisting}
\section*{Total Token Usage}
\textbf{Input tokens: }428
\\\textbf{Output tokens: }660
\\\textbf{Reasoning tokens: }512
\\\textbf{Sum of `total tokens': }1088
\section*{Experiment Timings}
\textbf{Overall Experiment} started at 1780486425389, ended at 1780486443180, lasting 17791ms (17.79 seconds)
\\\textbf{Iteration \#1} started at 1780486425390, ended at 1780486443180, lasting 17790ms (17.79 seconds)
\end{document}
